%! Introduction to Eigenvalues — Unit 6, Lesson 6.1 — Group Activity (Answer Key)
\documentclass[10pt]{article}
\usepackage{linalg-article}
\usepackage{linalg-key}   % pulls in -boxes; do NOT also load -boxes
\newcommand{\TallMath}[1]{$\displaystyle #1\rule[-1.4em]{0pt}{3.2em}$}
\newcommand{\xx}{\mathbf{x}}
\newcommand{\zero}{\mathbf{0}}

\begin{document}

\pageheader{Unit 6, Lesson 6.1}{Group Activity --- Answer Key}
\namepartnerperiod

\vspace{0.10in}

\begin{headlinebox}{blush}
\small\textbf{Finding the Special Directions --- No Guessing.} The method is always two steps.
\textbf{Step 1:} solve the characteristic equation $\det(A-\lambda I)=0$ for the \textbf{eigenvalues}
$\lambda$. \textbf{Step 2:} for each $\lambda$, solve $(A-\lambda I)\xx=\zero$ for an
\textbf{eigenvector}. Check with trace ($=\lambda_1+\lambda_2$) and determinant ($=\lambda_1\lambda_2$).
Show every step with your partner.
\end{headlinebox}

\vspace{0.08in}

\begin{tcolorbox}[colback=white, colframe=black!40, title=\textbf{Tier R --- Remediate},
    fonttitle=\bfseries, arc=1.5mm, left=3mm, right=3mm, top=2mm, bottom=2mm, breakable]
Warm up each half-skill.
\begin{enumerate}[leftmargin=1.7em, itemsep=9pt]
  \item \textbf{Diagonal matrix.} For $A=\begin{bsmallmatrix} 6 & 0\\ 0 & 2 \end{bsmallmatrix}$,
        $\det(A-\lambda I)=(6-\lambda)(2-\lambda)$. The eigenvalues are $\lambda=\ans{$6$}$ and
        $\lambda=\ans{$2$}$. An eigenvector for each is $\begin{bsmallmatrix}{\color{keyred}\mathbf{1}}\\[1pt]{\color{keyred}\mathbf{0}}\end{bsmallmatrix}$
        and $\begin{bsmallmatrix}{\color{keyred}\mathbf{0}}\\[1pt]{\color{keyred}\mathbf{1}}\end{bsmallmatrix}$. \emph{What do you notice about the diagonal?} \ansline{The eigenvalues \emph{are} the diagonal entries; the eigenvectors are the axes.}
  \item \textbf{Full method (symmetric).} Let $A=\begin{bsmallmatrix} 4 & 1\\ 1 & 4 \end{bsmallmatrix}$.
        Write $A-\lambda I$, compute $\det(A-\lambda I)$, and solve for both eigenvalues.
        \ansline{$A-\lambda I=\begin{bsmallmatrix}4-\lambda&1\\1&4-\lambda\end{bsmallmatrix}$; $\det=(4-\lambda)^2-1=\lambda^2-8\lambda+15=(\lambda-3)(\lambda-5)$, so $\lambda=3,5$.}
  \item Using your $\lambda$'s from problem~2, find an eigenvector for the \emph{larger} $\lambda$ by
        solving $(A-\lambda I)\xx=\zero$. \ansline{$\lambda=5$: $A-5I=\begin{bsmallmatrix}-1&1\\1&-1\end{bsmallmatrix}\Rightarrow x_1=x_2\Rightarrow\xx=\begin{bsmallmatrix}1\\1\end{bsmallmatrix}$.}
\end{enumerate}
\end{tcolorbox}

\begin{tcolorbox}[colback=white, colframe=black!40, title=\textbf{Tier A --- Approaching Proficiency},
    fonttitle=\bfseries, arc=1.5mm, left=3mm, right=3mm, top=2mm, bottom=2mm, breakable]
Run the full method, including the check.
\begin{enumerate}[leftmargin=1.7em, itemsep=9pt]
  \item \textbf{The whole job.} Let $A=\begin{bsmallmatrix} 4 & 2\\ 1 & 3 \end{bsmallmatrix}$.
        \begin{enumerate}[label=(\alph*), leftmargin=1.7em, itemsep=3pt]
          \item Find both eigenvalues from $\det(A-\lambda I)=0$.
          \item Find an eigenvector for each $\lambda$.
          \item \emph{Check:} does the sum of your $\lambda$'s equal the trace, and the product equal $\det A$?
        \end{enumerate}
        \ansline{(a) $\det(A-\lambda I)=(4-\lambda)(3-\lambda)-2=\lambda^2-7\lambda+10=(\lambda-2)(\lambda-5)\Rightarrow\lambda=2,5$.}
        \ansline{(b) $\lambda=5$: $A-5I=\begin{bsmallmatrix}-1&2\\1&-2\end{bsmallmatrix}\Rightarrow x_1=2x_2\Rightarrow\begin{bsmallmatrix}2\\1\end{bsmallmatrix}$; \quad $\lambda=2$: $A-2I=\begin{bsmallmatrix}2&2\\1&1\end{bsmallmatrix}\Rightarrow x_1=-x_2\Rightarrow\begin{bsmallmatrix}1\\-1\end{bsmallmatrix}$.}
        \ansline{(c) sum $2+5=7=$ trace $(4+3)$; product $2\cdot5=10=\det A\ (12-2)$. Both match.}
  \item \textbf{A $\lambda=0$ matrix (recall §5).} Let $A=\begin{bsmallmatrix} 2 & 4\\ 1 & 2 \end{bsmallmatrix}$.
        Find its eigenvalues. One of them is $0$ --- find the eigenvector for $\lambda=0$ (that is the
        direction $A$ crushes). What does $\lambda=0$ tell you about $\det A$ and invertibility?
        \ansline{$\det(A-\lambda I)=(2-\lambda)^2-4=\lambda^2-4\lambda=\lambda(\lambda-4)\Rightarrow\lambda=0,4$. For $\lambda=0$: $A\xx=\zero\Rightarrow x_1+2x_2=0\Rightarrow\begin{bsmallmatrix}2\\-1\end{bsmallmatrix}$. $\lambda=0$ means $\det A=0$, so $A$ is \emph{not invertible} --- a nonzero direction is crushed to $\zero$ (the §5 collapse).}
\end{enumerate}
\end{tcolorbox}

\begin{tcolorbox}[colback=white, colframe=black!40, title=\textbf{Tier E --- Extension},
    fonttitle=\bfseries, arc=1.5mm, left=3mm, right=3mm, top=2mm, bottom=2mm, breakable]
\textbf{When there are \emph{no} special directions.} A $90^\circ$ rotation is
$A=\begin{bsmallmatrix} 0 & -1\\ 1 & 0 \end{bsmallmatrix}$ (Lesson~5.3). It turns \emph{every} vector.
\begin{enumerate}[leftmargin=1.7em, itemsep=9pt]
  \item Write the characteristic equation $\det(A-\lambda I)=0$ and simplify it.
        \ansline{$\det\begin{bsmallmatrix}-\lambda&-1\\1&-\lambda\end{bsmallmatrix}=(-\lambda)(-\lambda)-(-1)(1)=\lambda^2+1=0$.}
  \item Try to solve it for a real number $\lambda$. What goes wrong?
        \ansline{$\lambda^2=-1$ has \emph{no real solution} (no real number squares to $-1$). So the matrix has no \emph{real} eigenvalues.}
  \item Explain \emph{geometrically} why this matrix has no real eigenvector --- and connect that to
        what a $90^\circ$ rotation does to arrows in the plane.
        \ansline{A $90^\circ$ rotation turns \emph{every} vector to a perpendicular direction, so no nonzero vector stays on its own line --- there is no direction that is only stretched. No real eigenvector, matching the ``no real $\lambda$'' from the equation.}
\end{enumerate}
\end{tcolorbox}

\begin{teachernote}
Tier~R splits the method into halves: read eigenvalues off a diagonal (item~1), find eigenvalues by the
characteristic equation (item~2), then find one eigenvector (item~3). Tier~A is the full job plus the
\emph{trace/determinant check}, and revisits the §5 callback with a $\lambda=0$ (singular) matrix
--- the eigenvector for $\lambda=0$ is exactly the nullspace direction. Tier~E is the conceptual
capstone: a rotation has no \emph{real} eigenvalues because it turns every vector (a natural place to
mention that complex eigenvalues $\pm i$ exist, though we stay with real ones). Common errors: sign
slips in $A-\lambda I$; forgetting the ``$-bc$'' term in the determinant; and, in A2, not recognizing
$\lambda(\lambda-4)$ gives $\lambda=0$. If time is short, Tier~A item~1 is the must-do.
\end{teachernote}

\end{document}
